% Master's laboratory report assessment template
% Keep the original review.tex preamble from \documentclass through
% \usepackage{siunitx}. Replace EVERYTHING after that package line
% with this file. Do not retain the old custom commands or document body.
% Suggested assessment weights below are editable, not institutional rules.

\definecolor{feedbackblue}{RGB}{0,75,125}
\definecolor{feedbackred}{RGB}{135,35,25}

\newcommand{\placeholder}[1]{\textit{[#1]}}
\newcommand{\strength}[1]{%
  \par\noindent\textbf{Strength:} #1\par\smallskip}
\newcommand{\observation}[1]{%
  \par\noindent\textbf{Observation:} #1\par\smallskip}
\newcommand{\requiredrevision}[1]{%
  \par\noindent\textbf{\textcolor{feedbackred}{Required revision:}} #1\par\smallskip}
\newcommand{\suggestion}[1]{%
  \par\noindent\textbf{\textcolor{feedbackblue}{Suggested improvement:}} #1\par\smallskip}
\newcommand{\studentresponse}[1]{%
  \par\noindent\textbf{Student response and location of changes:} #1\par\smallskip}
\newcommand{\revisioncheck}[1]{%
  \par\noindent\textbf{Assessor's revision check:} #1\par\medskip}

\newcounter{mainfeedback}
\newcommand{\mainfeedback}[1]{%
  \refstepcounter{mainfeedback}%
  \subsection*{M\arabic{mainfeedback}. #1}}
\newcounter{specificfeedback}
\newcommand{\specificfeedback}[2]{%
  \refstepcounter{specificfeedback}%
  \par\medskip\noindent\textbf{D\arabic{specificfeedback}. #1}\par
  #2\par\smallskip}

\title{Master's Laboratory Report\\Assessment and Revision Feedback}
\author{}
\date{\today}

\begin{document}
\maketitle
\vspace{-2em}

\section*{Report information}
\begin{description}[style=nextline,leftmargin=0pt]
  \item[Course / module] \placeholder{Course or module title}
  \item[Experiment / report title] \placeholder{Experiment and report title}
  \item[Student(s) / group] \placeholder{Names; student IDs if required; group}
  \item[Assessor] \placeholder{Name}
  \item[Submission] \placeholder{Date and report version}
  \item[Assessment round] \placeholder{Initial submission / revised submission}
\end{description}

\section*{Overall assessment}
\placeholder{Briefly assess how well the report addresses the experimental
objectives and demonstrates Master's-level scientific understanding.
Summarise the reliability of the results, quality of interpretation,
and the most important strengths and limitations. Assess achievement
of the learning objectives rather than novelty suitable for publication.}

\strength{\placeholder{Identify a specific scientific or methodological strength.}}
\strength{\placeholder{Identify a second strength, where appropriate.}}

\noindent\textbf{Main priorities for improvement:}
\begin{enumerate}[label=\arabic*.]
  \item \placeholder{Highest-priority scientific or methodological issue.}
  \item \placeholder{Second priority.}
  \item \placeholder{Third priority, if applicable.}
\end{enumerate}

\section*{Assessment criteria and points}
\noindent\textit{The categories and weights below are suggested for a reusable
laboratory-report rubric; they are not taken from the supplied grade scale.
Adapt them to the module's learning outcomes and communicate them before
assessment. Award points within each category and justify deductions in the
comments.}

\begin{center}
\renewcommand{\arraystretch}{1.25}
\begin{tabular}{@{}p{0.69\linewidth}rr@{}}
\hline
\textbf{Criterion} & \textbf{Maximum} & \textbf{Awarded} \\
\hline
Scientific background, objectives, and hypotheses & 15 & \placeholder{--} \\
Experimental methods and reproducibility & 20 & \placeholder{--} \\
Data processing and presentation of results & 20 & \placeholder{--} \\
Uncertainty analysis and critical quality control & 15 & \placeholder{--} \\
Discussion, physical interpretation, and conclusions & 20 & \placeholder{--} \\
Structure, scientific language, figures, and references & 10 & \placeholder{--} \\
\hline
\textbf{Total} & \textbf{100} & \placeholder{--} \\
\hline
\end{tabular}
\end{center}

\subsection*{Criterion-specific justification}
\begin{description}[style=nextline,leftmargin=0pt]
  \item[Scientific background, objectives, and hypotheses]
  \placeholder{Assess relevant theory, definitions, assumptions, and clarity
  of objectives; refer to main comments where appropriate.}
  \item[Experimental methods and reproducibility]
  \placeholder{Assess apparatus, calibration, measurement procedure, sample
  description, experimental conditions, and sufficient detail to reproduce
  the work. Assess safety documentation where relevant to the assignment.}
  \item[Data processing and presentation of results]
  \placeholder{Assess processing steps, equations, units, numerical accuracy,
  traceability of derived quantities, and effective tables and figures.}
  \item[Uncertainty analysis and critical quality control]
  \placeholder{Assess uncertainty estimates and propagation, repeatability,
  systematic effects, outlier treatment, and limitations of the measurements.}
  \item[Discussion, physical interpretation, and conclusions]
  \placeholder{Assess the physical reasoning, comparison with theory or
  reference data, interpretation of discrepancies, limitations, and whether
  the conclusions are supported by the results.}
  \item[Structure, scientific language, figures, and references]
  \placeholder{Assess organisation, clarity, notation, captions, axis labels,
  citation practice, and attribution of external material. Avoid counting
  the same deficiency again under another criterion without justification.}
\end{description}

\section*{Percentage-to-grade scale}
\noindent The following conversion reproduces the thresholds and evaluation
categories in the supplied HSBO grading sheet. The percentage is the points
awarded divided by the available points, multiplied by 100. With the
100-point rubric above, the total points equal the percentage.

\begin{center}
\renewcommand{\arraystretch}{1.15}
\begin{tabular}{@{}lll@{}}
\hline
\textbf{Percentage} & \textbf{Grade} & \textbf{Evaluation} \\
\hline
\(95 \leq p \leq 100\) & 1.0 & Very good (sehr gut) \\
\(90 \leq p < 95\) & 1.3 & Very good (sehr gut) \\
\(85 \leq p < 90\) & 1.7 & Good (gut) \\
\(80 \leq p < 85\) & 2.0 & Good (gut) \\
\(75 \leq p < 80\) & 2.3 & Good (gut) \\
\(70 \leq p < 75\) & 2.7 & Satisfactory (befriedigend) \\
\(65 \leq p < 70\) & 3.0 & Satisfactory (befriedigend) \\
\(60 \leq p < 65\) & 3.3 & Satisfactory (befriedigend) \\
\(55 \leq p < 60\) & 3.7 & Sufficient (ausreichend) \\
\(50 \leq p < 55\) & 4.0 & Sufficient (ausreichend) \\
\(0 \leq p < 50\) & 5.0 & Insufficient (nicht ausreichend) \\
\hline
\end{tabular}
\end{center}

\noindent\textbf{Rounding convention:}
\placeholder{State the applicable module rule. The supplied grading sheet
specifies thresholds but does not specify a rounding convention.}

\section*{Main comments}
\noindent\textit{Use one numbered comment per substantive issue. Explain the
issue, why it matters, and what constitutes an adequate revision. Clearly
distinguish required corrections from optional improvements. Duplicate the
comment blocks below as needed.}

\mainfeedback{\placeholder{Scientific understanding or interpretation}}
\observation{\placeholder{Describe the issue, cite the relevant report
location, and explain its scientific significance.}}
\requiredrevision{\placeholder{State a concrete, achievable correction and
what evidence or explanation is needed.}}
\studentresponse{\placeholder{Explain the change and give the page, section,
figure, equation, or line in the revised report; justify any disagreement.}}
\revisioncheck{\placeholder{Resolved / partly resolved / unresolved;
brief justification and any remaining action.}}

\mainfeedback{\placeholder{Methods, data processing, or uncertainty}}
\observation{\placeholder{Identify the affected measurement, calculation,
assumption, or uncertainty estimate and explain its effect on the results.}}
\requiredrevision{\placeholder{Specify the required clarification,
recalculation, uncertainty analysis, or limitation statement.}}
\studentresponse{\placeholder{Response and location of changes.}}
\revisioncheck{\placeholder{Resolution status and assessor's comment.}}

\mainfeedback{\placeholder{Further improvement}}
\observation{\placeholder{Describe a further issue where applicable.}}
\suggestion{\placeholder{Explain an optional improvement that would make the
report clearer or more informative.}}
\studentresponse{\placeholder{Response, if applicable.}}
\revisioncheck{\placeholder{Assessor's comment, if applicable.}}

\section*{Detailed comments}
\noindent\textit{Reference the assessed version explicitly. Use page,
section, figure, table, equation, or line numbers. Duplicate each block
as needed; numbering is automatic.}

\specificfeedback{Page \placeholder{X}, Section \placeholder{Y}}{%
  \observation{\placeholder{Specific issue in the text or reasoning.}}
  \requiredrevision{\placeholder{Required correction.}}
  \studentresponse{\placeholder{Response and revised location.}}
  \revisioncheck{\placeholder{Resolution status.}}
}

\specificfeedback{Figure \placeholder{X}, page \placeholder{Y}}{%
  \observation{\placeholder{Specific issue with the plotted data, units,
  uncertainty information, legend, or caption.}}
  \requiredrevision{\placeholder{Required correction.}}
  \studentresponse{\placeholder{Response and revised location.}}
  \revisioncheck{\placeholder{Resolution status.}}
}

\specificfeedback{Equation \placeholder{X}, page \placeholder{Y}}{%
  \observation{\placeholder{Specific issue with a definition, assumption,
  equation, notation, dimensional consistency, or numerical substitution.}}
  \suggestion{\placeholder{Suggested improvement; change to
  required revision if the issue affects scientific correctness.}}
  \studentresponse{\placeholder{Response and revised location.}}
  \revisioncheck{\placeholder{Assessor's comment.}}
}

\section*{Revision arrangements}
\begin{description}[style=nextline,leftmargin=0pt]
  \item[Assessment status]
  \placeholder{Final assessment / provisional assessment pending revision /
  formative feedback without a grade.}
  \item[Revision opportunity and deadline]
  \placeholder{State whether resubmission is available, the deadline,
  and any restrictions under the module rules.}
  \item[Required submission material]
  \placeholder{For example: revised report and point-by-point responses
  referencing comment IDs M1, M2, D1, etc.; indicate whether a marked-up
  version or processing files are required.}
  \item[Scope of reassessment]
  \placeholder{State whether the entire report or only specified criteria
  will be reassessed and whether any grade cap applies. Do not assume
  that revision automatically changes the grade.}
\end{description}

\section*{Assessment record}
\begin{center}
\renewcommand{\arraystretch}{1.3}
\begin{tabular}{@{}p{0.40\linewidth}p{0.25\linewidth}p{0.25\linewidth}@{}}
\hline
 & \textbf{Initial submission} & \textbf{Revised submission} \\
\hline
Assessment date & \placeholder{Date} & \placeholder{Date / N/A} \\
Total points / available points & \placeholder{-- / 100} & \placeholder{-- / 100} \\
Percentage & \placeholder{--\%} & \placeholder{--\%} \\
Grade & \placeholder{Grade / provisional} & \placeholder{Grade / N/A} \\
\hline
\end{tabular}
\end{center}

\noindent\textbf{Final recorded grade:} \placeholder{Grade / not yet assigned}

\medskip
\noindent\textbf{Final assessor's comment:}
\placeholder{Record the extent to which the required revisions were
addressed and justify the final assessment. If relevant, identify one or
two priorities for future laboratory work.}

\medskip
\noindent\textbf{Assessor / date:} \placeholder{Name and date}

\end{document}
